\chapter*{Abstract}
\addcontentsline{toc}{chapter}{Abstract}

TrustFin AI is a completed loan decision-support system that combines product-specific credit prediction with post-hoc explanations and quantitative checks on those explanations. The project addresses a practical limitation of using complex models for lending: a prediction can be useful for ranking risk while remaining difficult for an applicant or reviewing officer to interpret. A single explanation is also not sufficient evidence that the explanation is faithful, stable, consistent across methods, or insensitive to a protected characteristic.

The implemented system supports home-loan and personal-loan applications. Each product uses a separately trained XGBoost classifier and a product-specific preprocessing pipeline. Derived repayment features include monthly instalment and fixed-obligation-to-income ratio (FOIR); the home-loan pipeline also derives loan-to-value ratio. The prediction service returns the model decision and class probability, and produces local explanations using TreeSHAP and tabular LIME. An explanation-validation stage calculates faithfulness, stability under a small seeded perturbation, top-feature consistency between SHAP and LIME, and an individual counterfactual gender-flip score. These values are combined using documented weights to form a composite Trust Score. The result, selected applicant fields, explanation data, and trust measures can be stored through SQLAlchemy using local SQLite or a configured PostgreSQL database. A React dashboard presents results to applicant and bank-officer views and the backend can generate a PDF report.

The project datasets are generated by repository scripts and contain 3,000 records each; they are not records from a lending institution. On the stored stratified test evaluations, the home-loan model reports 97.00\% accuracy and 0.9978 ROC-AUC, while the personal-loan model reports 97.33\% accuracy and 0.9963 ROC-AUC. These values describe performance on generated data and do not establish performance on real applicants. The composite score is an engineering indicator defined by the project, not an independently validated measure of safety or lending quality. The report therefore distinguishes predictive performance, prediction confidence proxy, explanation checks, and real-world fairness evidence.

\noindent\textbf{Keywords:} Explainable AI, credit scoring, XGBoost, SHAP, LIME, explanation validation, fairness, Trust Score.

\clearpage
